% II. MARCO TEÓRICO. Lo escribe la skill tesis-marco.

\section{Antecedentes de la investigación}
\subsection{Antecedentes internacionales}
% Mínimo 5, de los últimos meta.antecedentes_ventana_anios años. Un párrafo por antecedente,
% \textcite{clave}, parafraseado y sin página: autor y año -> lugar -> objetivo -> diseño,
% población/muestra -> instrumentos -> resultados con cifras -> conclusión vinculada.

\subsection{Antecedentes nacionales}
% Mínimo 3; tesis o artículos peruanos recientes. Misma plantilla.

\section{Bases teóricas}
% Un \subsection por variable (VI y VD) con sus dimensiones e indicadores y fórmulas.
% Mínimo 3 definiciones por variable de autores distintos, integradas.
%
% Figura APA (el \caption va ANTES de la imagen):
% \begin{figure}[H]
%   \caption{Arquitectura del sistema propuesto}
%   \label{fig:arquitectura}
%   \centering\includegraphics[width=0.85\textwidth]{arquitectura}
%   \nota{Elaboración propia.}
% \end{figure}
%
% Tabla APA (solo líneas horizontales):
% \begin{table}[H]
%   \caption{Comparación de modelos de lenguaje}
%   \label{tab:modelos}
%   \begin{tabularx}{\textwidth}{@{}lYr@{}}
%     \toprule
%     \textbf{Modelo} & \textbf{Descripción} & \textbf{Costo} \\
%     \midrule
%     A & ... & 0.00 \\
%     \bottomrule
%   \end{tabularx}
%   \nota{Adaptado de \textcite{clave}.}
% \end{table}

\section{Definición de términos básicos}
% 10 a 20 términos en orden alfabético, cada uno citado:
% \textbf{Término.} Definición parafraseada \parencite{clave}.
